\chapter{Reference sheets}
\chaptertag{Definitions, algorithms, and proof moves}

\section*{Core notation}

\begin{center}
\begin{tabularx}{\textwidth}{p{.22\textwidth}X}
\toprule
Notation & Meaning \\
\midrule
$a\divides b$ & $b=ak$ for some integer $k$ \\
$\gcd(a,b)$ & greatest positive common divisor \\
$a\congr b\pmod n$ & $n\divides(a-b)$ \\
$(\Z/n\Z)^\times$ & invertible residue classes modulo $n$ \\
$\varphi(n)$ & number of units modulo $n$ \\
$v_p(m)$ & exponent of $p$ in $m$ \\
$\pi(x)$ & number of primes at most $x$ \\
$\vartheta(x)$ & $\sum_{p\le x}\log p$ \\
$\psi(x)$ & $\sum_{p^k\le x}\log p$ \\
$f(x)\asymp g(x)$ & their ratio stays between positive constants eventually \\
\bottomrule
\end{tabularx}
\end{center}

\section*{Algorithms at a glance}

\textbf{Extended Euclid.} Divide until zero; the last nonzero remainder is the gcd; substitute backward to obtain $ax+by=\gcd(a,b)$.

\textbf{Sieve of Eratosthenes.} To find the primes through $N$, repeatedly
circle the least uncrossed integer and cross out its larger multiples. Stop
after the sieving primes $p\le\sqrt N$; a composite at most $N$ must have a
prime factor in that range.

\textbf{Existence versus a witness.} The classical dichotomy for
$x=(\sqrt2)^{\sqrt2}$ proves that one of
$(\sqrt2,\sqrt2)$ and $(x,\sqrt2)$ has a rational power, but does not say
which. The explicit choice
$a=\sqrt2$, $b=\log_{\sqrt2}3$ gives $a^b=3$, and unique factorization proves
that $b$ is irrational.

\textbf{Continued fraction.} Repeatedly take the integer part, subtract it, and invert. Convergents satisfy
\[
p_k=a_kp_{k-1}+p_{k-2},\qquad q_k=a_kq_{k-1}+q_{k-2}.
\]
The recurrences are the columns of the matrix product
\[
\begin{pmatrix}a_0&1\\1&0\end{pmatrix}\cdots
\begin{pmatrix}a_k&1\\1&0\end{pmatrix}
=\begin{pmatrix}p_k&p_{k-1}\\q_k&q_{k-1}\end{pmatrix},
\]
whose determinant is $(-1)^{k+1}$. Each factor lies in
$\operatorname{GL}_2(\Z)$; a product with an even number of factors lies in
$\operatorname{SL}_2(\Z)$.
For $A/B=[a_0;\dots,a_n]$, the Euclidean remainders obey
\[
r_k=(-1)^k(q_kA-p_kB).
\]
For irrational $x$, the zero-indexed convergents form alternating monotone bounds:
\[
c_0<c_2<c_4<\cdots<x<\cdots<c_5<c_3<c_1.
\]
Legendre's test says that a reduced fraction satisfying
$\abs{x-p/q}<1/(2q^2)$ must be a convergent. Lagrange's theorem says that a
simple continued fraction is eventually periodic exactly for a quadratic
irrational; Galois says it is purely periodic exactly when $x>1$ and
$-1<x'<0$.

\textbf{Fixed points and prefixes.} For $\det M\ne0$,
$T_M(t)=(at+b)/(ct+d)$ satisfies $T_M\circ T_N=T_{MN}$.
An irrational fixed point of a nonscalar rational $M$ obeys the genuine
quadratic $ct^2+(d-a)t-b=0$. If $x=T_N(x)$ and $y=T_M(x)$, then
\[
 y=T_{MNM^{-1}}(y).
\]
A finite continued-fraction prefix therefore preserves quadratic
irrationality by conjugating the period matrix.

\textbf{Reduced quadratic states.} Write $x_n=(P_n+\sqrt D)/Q_n$ with
$Q_n\divides D-P_n^2$. With $a_n=\floor{x_n}$, the exact update is
\[
 P_{n+1}=a_nQ_n-P_n,\qquad Q_{n+1}=(D-P_{n+1}^2)/Q_n.
\]
Every quadratic irrational eventually has a reduced tail.
For reduced states, $|P_n|<\sqrt D$ and $0<Q_n<2\sqrt D$,
so only finitely many states occur. A reduced successor $z$ determines
its reduced predecessor:
$a=\floor{-1/z'}$ and $u=a+1/z$.
Finiteness gives eventual repetition; unique predecessors make a
reduced starting state purely periodic.

\textbf{Typical information rate.} For almost every $x$, the Khinchin--L\'evy
law gives
\[
 \lim_{n\to\infty}q_n^{1/n}=e^{\pi^2/(12\log2)}\approx3.275823,
 \qquad
 \lim_{n\to\infty}\abs{x-p_n/q_n}^{-1/n}
 =e^{\pi^2/(6\log2)}\approx10.731.
\]
So a continued-fraction term earns about $1.0306$ decimal digits on average.
For the first $100{,}000$ partial quotients of $\pi$, computation gives
$G_n(\pi)\approx2.67956$ and $q_n^{2/n}\approx10.67835$; neither limiting law
has been proved for $\pi$. In contrast,
$e=[2;1,2,1,1,4,1,1,6,\dots]$ has
$G_{3m}(e)\sim(2m/e)^{1/3}\to\infty$ and $q_{3m}^{2/(3m)}\to\infty$.

\textbf{Gosper's streaming arithmetic.} Represent a unary rational
transformation by $M=\begin{psmallmatrix}A&B\\C&D\end{psmallmatrix}$.  Reading
an input digit $m$ and emitting a certified output digit $k$ update the state by
\[
 M\longmapsto M\begin{pmatrix}m&1\\1&0\end{pmatrix},
 \qquad
 M\longmapsto\begin{pmatrix}0&1\\1&-k\end{pmatrix}M.
\]
For two inputs use
$(axy+bx+cy+d)/(exy+fx+gy+h)$; its eight integer coefficients are closed
under the same absorb-or-emit operations.  Exact interval bounds certify a
digit when every possible value has the same floor.  Interval-valued output is
the safe fallback at integer boundaries where a digit-only process can stall.

\textbf{Calendar design.} If a tropical year is $365+\alpha$ days, a cycle
with $p$ leap days in $q$ years asks for $p/q\approx\alpha$. For
$\alpha\approx0.2421988$, the convergents begin
\[
 \frac14,\quad\frac7{29},\quad\frac8{33},\quad\frac{31}{128}.
\]
The Julian calendar uses $1/4$. Requiring $q=100Q$ means approximating
$100\alpha$; its convergent $97/4$ yields the Gregorian rate $97/400$.

\textbf{Planetary lapping time.} If an inner planet takes $V$ days for one
Sun-centered lap and Earth takes $E>V$ days, the waiting time $S$ between
successive alignments (the \emph{synodic period}) satisfies
\[
  \frac1S=\frac1V-\frac1E,
  \qquad \frac SE=T(E/V),\qquad T(r)=\frac1{r-1}.
\]
For Earth and Venus, $E/V\approx13/8$ and $S/E\approx8/5$; the five
successive conjunction directions advance by about $3/5$ of a turn and draw a
near-pentagram.  The Dresden Codex uses the exact nominal calendar identities
\[
  5\cdot584=8\cdot365=2920,
  \qquad
  65\cdot584=104\cdot365=146\cdot260=37960.
\]

\textbf{Square-root periods.} For a nonsquare integer $D\ge2$,
\[
 \sqrt D=[a_0;\overline{a_1,\ldots,a_{L-1},2a_0}],
 \qquad a_j=a_{L-j}\quad(1\le j<L),
\]
where $L$ is the least period length. The block before $2a_0$ is a
palindrome. The complete quotients satisfy
$Q_n=1$ exactly when $L\divides n$, for $n\ge1$.

\textbf{Pell bridge.} If $p^2-Dq^2=\pm1$ with $p,q>0$, then
\[
 \abs{\sqrt D-p/q}=\frac{1}{q^2(\sqrt D+p/q)}<\frac1{2q^2},
\]
so Legendre makes $p/q$ a convergent. The exact identity
\[
 p_{n-1}^2-Dq_{n-1}^2=(-1)^nQ_n
\]
shows that the Pell convergents have indices $mL-1$ and sign $(-1)^{mL}$.
The first $+1$ solution is at $L-1$ for even $L$ and $2L-1$ for odd $L$;
the $-1$ equation is solvable exactly when $L$ is odd.
Stop before the final $2a_0$. If $(u,v)$ is the first $+1$ solution,
every positive $+1$ solution is a power of $u+v\sqrt D$, or equivalently
is generated by $\begin{psmallmatrix}u&Dv\\v&u\end{psmallmatrix}$.

\textbf{Cattle reduction.} Archimedes' seven proportions give the total
$50\,389\,082t$. The square condition forces $t=4\,456\,749r^2$;
the triangular condition is equivalent to
\[
 X^2-4\,729\,494Z^2=1,\qquad Z=9314r.
\]
The divisibility condition $9314\divides Z$ must be retained. The first
admissible power of the fundamental unit is $2329$; the resulting
smallest total has $206\,545$ digits. See
Section~\ref{sec:archimedes-cattle} and Vardi's exposition.

\textbf{Base-$b$ expansion.} Long division evolves by
$br_n=qd_{n+1}+r_{n+1}$ with $0\le r_{n+1}<q$. Repeated remainders make every
rational expansion eventually periodic. A reduced fraction $m/q$ is purely
periodic exactly when $\gcd(q,b)=1$; the reverse direction follows from
Euler's theorem.

\textbf{Linear congruence.} For $ax\congr b\pmod n$, first compute $d=\gcd(a,n)$. There are solutions iff $d\divides b$. Divide by $d$, reduce the modulus to $n/d$, and invert the new coefficient.

\textbf{CRT.} With $N=\prod n_i$, take $N_i=N/n_i$ and choose $M_i=N_i^{-1}\pmod{n_i}$. Then $x=\sum a_iN_iM_i\pmod N$.

\section*{Proof moves worth naming}

\begin{itemize}
  \item \textbf{Minimal counterexample / strong induction}: factor a composite object into smaller ones.
  \item \textbf{B\'ezout insertion}: replace $1$ by $ax+by$ to transfer divisibility.
  \item \textbf{Exponent comparison}: use unique factorization or $v_p$ to compare products and quotients.
  \item \textbf{M\"obius monotonicity}: compute the derivative of $(at+b)/(ct+d)$ to determine how a continued-fraction tail moves its value.
  \item \textbf{Conjugation}: transport a fixed-point equation through a change of variable using $MNM^{-1}$.
  \item \textbf{Finite states with reversible steps}: bound integer states to force eventual repetition, then recover predecessors to remove any initial segment.
  \item \textbf{Equal-block partition}: cosets and rotation orbits turn a count into divisibility.
  \item \textbf{Count two ways}: binomial identities become transparent when both sides count the same set.
  \item \textbf{Logarithmic weighting}: products of primes become the sums $\vartheta$ and $\psi$.
  \item \textbf{Dyadic decomposition}: split a range into $(x/2,x],(x/4,x/2],\dots$ and sum a geometric series.
\end{itemize}

\section*{The two Fermat proofs}

\begin{center}
\small
\begin{tabularx}{\textwidth}{p{.18\textwidth}XX}
\toprule
 & Necklace proof & Group proof \\
\midrule
Objects & $a^p$ labeled colorings & nonzero residues modulo $p$ \\
Structure & cyclic rotation orbits & group $(\Z/p\Z)^\times$ \\
Exceptional part & $a$ constant colorings & identity/powers inside cyclic subgroups \\
Equal blocks & nonconstant orbits have size $p$ & element orders divide $p-1$ \\
Conclusion & $a^p\congr a\pmod p$ & $a^{p-1}\congr1\pmod p$ for $p\ndivides a$ \\
\bottomrule
\end{tabularx}
\end{center}
